\documentclass{article}
\usepackage{amsmath}
\usepackage[hidelinks]{hyperref}
\title{Do not erase velocity at a stopped-robot cut}
\date{}
\begin{document}
\maketitle

\paragraph{Target and pins.}
Use the four departure requests from the published
\href{https://arxiv.org/html/2609.29939v1}{KTH managed-removal examples}.
At the request cut only, counterfactually stop the departing body at
its current pose and remove its radio. Retain every surviving
position and velocity from the unmodified point-mass run. The full
geometric graph already splits in two, so continuous global
connectivity is not demanded retroactively. Preserve all old surviving
established edges and the separation between all live bodies and the
stationary failed body throughout. Restore a cross-component link
at the source activation margin and settle all survivors at rest.
The declared speed cap is $V=1$ and acceleration caps are $1/2,1,2$
in simulation units. These are constraints of the counterfactual,
not claims about physical limits of the published vehicles. No noise,
disturbance, position or velocity error is imposed; this is a
full-state capability comparison with synchronous immediate control.
No further arbitrary collision hazards are considered.
The source has nonzero differences in velocities on old surviving
edges, so insisting on unchanged relative offsets from the cut
would itself violate differentiability. We constrain clearance and
links, not a motion-incompatible exact old shape.

\paragraph{Universal directional bound for settled service.}
Fix any candidate cross pair with initial distance $d>R$ and let
$e$ point from the first member to the second. Let the closing
velocity projections initially be $u_i=e\cdot v_i$ and
$u_j=-e\cdot v_j$. A still-connected rest endpoint at horizon $T$
requires the integrated sum of projected closing motion at least
$d-R$. Each such projected velocity is bounded above at real time
$t$ by
\[
\min\{V,\ u_i+At,\ A(T-t)\}.
\]
The terms follow from bounded speed, initial signed velocity, and
the same input bound backwards from zero terminal velocity. The
candidate set can satisfy service only if at least one cross pair
can close by $d-R$ under these upper integrals. All survivors must
also have enough time to stop. Relaxing a signed projection to the
full initial speed gives a weaker necessary test. No collision path
or fixed control policy is assumed in either lower bound. Beyond
the longest braking time all upper integrals are monotone in $T$;
the piecewise affine intersections integrate directly. This gives
a lower deadline bound in the stated continuous model, rather than
a failure found by numerical controller search.

\paragraph{Constructive checked upper bound.}
Each survivor first uses maximum braking along its own velocity
until stationary. Between the individual stopping times and
acceleration switches, each relative trajectory is quadratic.
The deviation of a relative segment from its chord is at most
$\|a_i-a_j\|\Delta^2/8$. Lower chord distances minus this bound
check clearance; maximum endpoint norm plus it checks every
surviving established edge. Pointwise speed follows from the
endpoint speed bound and affine velocity on each segment.
Next each component applies one synchronized bounded rest-to-rest
translation toward a closest cross pair in the braked configuration.
It suffices to move each component half the remaining gap.
The triangular speed profile switches once (or cruises at the speed
cap if necessary), and preserves all internal offsets during this
phase. The same pairwise continuous bounds check both motion
phases, including the static failed body. At the end there is a
bridge with positive slack relative to nominal sensing range.
This is not an optimization claim for the brake-then-translate
controller or a claim about a switched control law outside these
four initial poses.

\paragraph{Positive and negative matched deadlines.}
All twelve pose/cap cases produce a witness within 0.001 modeled
time units of the projected lower bound. For example at acceleration
cap one the planar request has a speed-only static estimate
of about 0.22960, a warm settled lower bound about 0.95879, and
a verified settled trajectory about 0.95929. The source reports
small but nonzero velocities. The large timing correction is
acceleration plus coming to rest, not the small braking prelude
alone. If a separate authored acceptance law requires an intact
rested link for 0.25 units and a matching callback, that dwell is
added after the verified stop; the replay records an infeasible and
feasible-by deadline in the same task for each pose/cap. A missing
callback still does not establish qualified service. There is no
declaration that the published robot uses that transport protocol.
The lower/upper numerical bins come with slack between the
required negative/positive deadlines and do not assert
last-digit platform-independent precision.

\paragraph{Relation to established controllers.}
The published manager prevents the abrupt mode by maintaining its
old support until bridge activation. The connected-start theorem in
\href{https://arxiv.org/pdf/2609.19912}{Geurts et al.} is a maintenance
baseline; it does not supply the disconnection recovery initial
condition. The physical deadline obstruction applies also to
any planner given the same bounded plant and terminal requirements.
\href{https://en.endo-lab.jp/research/detail/499}{The authors' account}
of Maeda, Endo and Matsuno's recovery uses local neighbor behavior
and the failed member's previous direction (and can wait).
The exact published crash-control implementation has not been
replayed here. The ability of local crews to choose or track this
full-state witness under delayed evidence and a real failure
transition remains open. No physical, independent-mission, complete
R6 or field claim follows.
\end{document}
